\documentclass[12pt,a4paper]{article}

% ============================================================
% PACKAGES
% ============================================================
\usepackage[utf8]{inputenc}
\usepackage[T1]{fontenc}
\usepackage[french]{babel}
\usepackage{graphicx}
\usepackage{amsmath,amssymb}
\usepackage{booktabs}
\usepackage{geometry}
\usepackage{hyperref}
\usepackage{xcolor}
\usepackage{float}
\usepackage{caption}
\usepackage{subcaption}
\usepackage{enumitem}
\usepackage{fancyhdr}
\usepackage{titlesec}
\usepackage{listings}
\usepackage{array}

\geometry{margin=2.5cm}

% ============================================================
% STYLING
% ============================================================
\hypersetup{
    colorlinks=true,
    linkcolor=blue!70!black,
    urlcolor=blue!70!black,
    citecolor=blue!70!black
}

\pagestyle{fancy}
\fancyhf{}
\rhead{Link Prediction — Actor Co-occurrence Network}
\lhead{}
\rfoot{Page \thepage}

\titleformat{\section}{\Large\bfseries\color{blue!70!black}}{}{0em}{}[\titlerule]
\titleformat{\subsection}{\large\bfseries\color{blue!50!black}}{}{0em}{}

\lstset{
    language=Python,
    basicstyle=\ttfamily\small,
    keywordstyle=\color{blue!70!black},
    commentstyle=\color{green!50!black},
    stringstyle=\color{red!60!black},
    frame=single,
    breaklines=true,
    backgroundcolor=\color{gray!5}
}

% ============================================================
% TITLE PAGE
% ============================================================
\begin{document}

\begin{titlepage}
\centering
\vspace*{3cm}
{\Huge\bfseries\color{blue!70!black} Prédiction de Liens\\[0.3cm] dans un Réseau d'Acteurs\par}
\vspace{1cm}
{\Large Compétition Kaggle — Link Prediction\par}
\vspace{2cm}
{\large\textit{Approche basée sur le Feature Engineering\\et le Machine Learning Classique}\par}
\vspace{3cm}
{\large\today\par}
\vfill
{\small Score Kaggle obtenu : \textbf{AUC = 0.82}\par}
\end{titlepage}

\tableofcontents
\newpage

% ============================================================
% 1. INTRODUCTION
% ============================================================
\section{Introduction}

\subsection{Contexte de la compétition}

Cette compétition Kaggle porte sur la \textbf{prédiction de liens manquants} dans un réseau de co-occurrence d'acteurs. Le graphe sous-jacent est construit à partir de Wikipedia : chaque nœud représente un acteur, et une arête entre deux nœuds indique que ces acteurs apparaissent sur la même page Wikipedia. Cette co-occurrence sert de proxy pour leur participation conjointe dans un film, une relation personnelle, ou leur niveau de notoriété.

Des arêtes ont été supprimées aléatoirement de ce graphe. L'objectif est de \textbf{reconstruire les liens manquants} en utilisant à la fois la structure du graphe et les attributs des nœuds.

\subsection{Données disponibles}

\begin{table}[H]
\centering
\begin{tabular}{lp{9cm}}
\toprule
\textbf{Fichier} & \textbf{Description} \\
\midrule
\texttt{node\_information.csv} & 3\,597 nœuds × 932 features binaires (présence de mots-clés Wikipedia) \\
\texttt{train.txt} & 10\,496 paires étiquetées (5\,248 positives + 5\,248 négatives) \\
\texttt{test.txt} & 3\,498 paires à prédire \\
\bottomrule
\end{tabular}
\caption{Description des fichiers de données}
\end{table}

\subsection{Métrique d'évaluation}

La métrique utilisée est l'\textbf{AUC-ROC} (Area Under the Receiver Operating Characteristic Curve). Cette métrique mesure la capacité du modèle à distinguer les paires positives (lien existant) des paires négatives (pas de lien), indépendamment du seuil de classification. On prédit donc des \textbf{probabilités} et non des classes binaires.

\subsection{Contraintes}

\begin{itemize}
    \item Pas de deep learning (GNN, réseaux de neurones, transformers)
    \item Focus sur le \textbf{machine learning classique} : gradient boosting, random forest, régression logistique
    \item Priorité au \textbf{feature engineering} plutôt qu'à la complexité du modèle
\end{itemize}

% ============================================================
% 2. ANALYSE EXPLORATOIRE
% ============================================================
\newpage
\section{Analyse Exploratoire des Données}

\subsection{Structure du graphe}

Le graphe construit à partir des arêtes positives du train présente les caractéristiques suivantes :

\begin{table}[H]
\centering
\begin{tabular}{lr}
\toprule
\textbf{Propriété} & \textbf{Valeur} \\
\midrule
Nombre de nœuds & 3\,597 \\
Nombre d'arêtes & 5\,248 \\
Densité & 0.000811 (très sparse) \\
Composantes connexes & 1 (graphe connexe) \\
Degré moyen & 2.92 \\
Degré médian & 2 \\
Degré maximum & 361 (nœud hub) \\
\bottomrule
\end{tabular}
\caption{Statistiques du graphe d'acteurs}
\end{table}

Le graphe est \textbf{extrêmement creux} (densité de 0.08\%) mais \textbf{connexe}. La distribution des degrés suit approximativement une \textbf{loi de puissance} (power law), caractéristique des réseaux sociaux : la majorité des acteurs ont très peu de connexions, tandis que quelques hubs (acteurs très connus) en ont des centaines.

\begin{figure}[H]
\centering
\includegraphics[width=\textwidth]{figures/fig_01.png}
\caption{Distribution des degrés du graphe. \textbf{Gauche} : histogramme montrant la concentration sur les petits degrés. \textbf{Centre} : représentation log-log confirmant un comportement de type loi de puissance. \textbf{Droite} : courbe de rang montrant la décroissance rapide.}
\label{fig:degree}
\end{figure}

\subsection{Analyse des features de nœuds}

Chaque acteur est caractérisé par un vecteur binaire de 932 dimensions, représentant la présence ou l'absence de mots-clés extraits de sa page Wikipedia.

\begin{itemize}
    \item \textbf{Densité} : 0.58\% — les vecteurs sont extrêmement creux
    \item \textbf{Mots-clés par acteur} : en moyenne 5.4 ($\sigma$ = 3.25)
    \item Les features sont strictement binaires : $\{0, 1\}$
\end{itemize}

\begin{figure}[H]
\centering
\includegraphics[width=\textwidth]{figures/fig_02.png}
\caption{Distribution des features de nœuds. \textbf{Gauche} : nombre de mots-clés par acteur. \textbf{Droite} : nombre d'acteurs par mot-clé (features non nulles uniquement).}
\label{fig:features}
\end{figure}

\subsection{Analyse du jeu d'entraînement}

Le jeu d'entraînement est \textbf{parfaitement équilibré} : 5\,248 paires positives et 5\,248 paires négatives. J'ai analysé les propriétés des nœuds selon leur label pour identifier des signaux discriminants.

\begin{figure}[H]
\centering
\includegraphics[width=\textwidth]{figures/fig_03.png}
\caption{Distribution des degrés des nœuds source et target par label. Les paires positives impliquent davantage de nœuds à haut degré, surtout côté target.}
\label{fig:labels}
\end{figure}

\textbf{Observation clé} : les paires positives (arêtes existantes) impliquent des nœuds de degré significativement plus élevé que les paires négatives. Le degré moyen du nœud target est de 31.9 pour les positives vs 2.9 pour les négatives — un ratio de 11:1. Cela confirme l'effet \textit{preferential attachment} : les nœuds fortement connectés attirent davantage de liens.

% ============================================================
% 3. FEATURE ENGINEERING
% ============================================================
\newpage
\section{Feature Engineering}

Le feature engineering constitue le \textbf{cœur du projet} et la principale source de performance. J'ai développé 36 features réparties en 5 familles complémentaires.

\subsection{Prévention du data leakage}

\textbf{Problème identifié} : les paires positives du train sont par définition des arêtes du graphe. Si on calcule les features sur le graphe complet, les features comme \texttt{shortest\_path\_length = 1} identifient trivialement les arêtes existantes, donnant une AUC de 1.0 en validation croisée, mais ne généralisant pas au test.

\textbf{Solution} : pour chaque paire d'entraînement $(u, v)$, je \textbf{retire temporairement l'arête} $(u, v)$ du graphe avant de calculer les features topologiques locales (voisins communs, plus court chemin, etc.), puis je la restaure. Cela force le modèle à apprendre des signaux indirects.

\begin{lstlisting}[caption={Principe du retrait d'arête pour éviter le leakage}]
# Pour chaque paire (s, t) du train :
if G.has_edge(s, t):
    G.remove_edge(s, t)      # Retrait temporaire
    features = compute(G, s, t)  # Calcul sans l'arete
    G.add_edge(s, t)          # Restauration
\end{lstlisting}

\subsection{Famille 1 : Features topologiques (12 features)}

Ces features exploitent la structure locale du graphe autour des deux nœuds.

\begin{table}[H]
\centering
\small
\begin{tabular}{p{4.5cm}p{6cm}p{3cm}}
\toprule
\textbf{Feature} & \textbf{Formule} & \textbf{Ratio pos/neg} \\
\midrule
\texttt{common\_neighbors} & $|N(u) \cap N(v)|$ & 10.04 \\
\texttt{jaccard\_coefficient} & $\frac{|N(u) \cap N(v)|}{|N(u) \cup N(v)|}$ & 2.70 \\
\texttt{adamic\_adar\_index} & $\sum_{w \in N(u) \cap N(v)} \frac{1}{\log(\deg(w))}$ & 18.97 \\
\texttt{resource\_allocation} & $\sum_{w \in N(u) \cap N(v)} \frac{1}{\deg(w)}$ & 35.72 \\
\texttt{preferential\_attachment} & $\deg(u) \times \deg(v)$ & 9.38 \\
\texttt{degree\_source/target} & Degrés individuels & 1.6 / 10.8 \\
\texttt{degree\_sum/diff/min/max} & Combinaisons de degrés & Variable \\
\texttt{shortest\_path\_length} & Distance dans le graphe (cutoff=6) & 0.87 \\
\bottomrule
\end{tabular}
\caption{Features topologiques. Le ratio pos/neg mesure le pouvoir discriminant.}
\end{table}

\textbf{Justification} :
\begin{itemize}
    \item \textbf{Common neighbors} : deux acteurs partageant des voisins sont très susceptibles d'être liés. C'est l'heuristique de base en link prediction.
    \item \textbf{Adamic-Adar} : pondère les voisins communs par l'inverse du log de leur degré. Un voisin rare (peu de connexions) est plus informatif qu'un hub.
    \item \textbf{Resource Allocation} : similaire à Adamic-Adar mais utilise $1/\deg(w)$ — pénalise encore plus les hubs.
    \item \textbf{Preferential Attachment} : capture l'effet ``les riches s'enrichissent'' — les acteurs populaires forment plus de liens entre eux.
    \item \textbf{Shortest path} : un chemin court suggère une proximité structurelle. Calculé avec une borne supérieure de 6 pour la performance.
\end{itemize}

\subsection{Famille 2 : Features d'attributs (9 features)}

Ces features mesurent la similarité entre les vecteurs de mots-clés Wikipedia des deux acteurs.

\begin{table}[H]
\centering
\small
\begin{tabular}{p{4.5cm}p{7cm}p{2cm}}
\toprule
\textbf{Feature} & \textbf{Description} & \textbf{Ratio} \\
\midrule
\texttt{cosine\_similarity} & $\frac{\mathbf{f}_u \cdot \mathbf{f}_v}{\|\mathbf{f}_u\| \|\mathbf{f}_v\|}$ & 0.98 \\
\texttt{common\_keywords} & $\sum_i f_u^{(i)} \times f_v^{(i)}$ & 0.99 \\
\texttt{union\_keywords} & $|\{i : f_u^{(i)} = 1 \text{ ou } f_v^{(i)} = 1\}|$ & 1.01 \\
\texttt{jaccard\_features} & common / union & 1.02 \\
\texttt{feature\_diff\_l2} & $\|\mathbf{f}_u - \mathbf{f}_v\|_2$ & 1.00 \\
\texttt{keyword\_count\_*} & Nombre de mots-clés par nœud & $\sim$1.0 \\
\bottomrule
\end{tabular}
\caption{Features basées sur les attributs de nœuds}
\end{table}

\textbf{Observation} : ces features ont un pouvoir discriminant faible individuellement (ratios proches de 1.0). Cependant, elles apportent un signal \textbf{complémentaire} à la topologie : deux acteurs du même univers cinématographique partagent des mots-clés similaires même sans connexion directe dans le graphe observé.

\subsection{Famille 3 : Features de communauté (3 features)}

J'utilise l'algorithme de \textbf{Louvain} pour détecter les communautés dans le graphe. Cet algorithme optimise la modularité et identifie des groupes de nœuds densément connectés.

\begin{table}[H]
\centering
\begin{tabular}{p{5cm}p{6cm}p{2cm}}
\toprule
\textbf{Feature} & \textbf{Description} & \textbf{Ratio} \\
\midrule
\texttt{same\_community} & 1 si même communauté, 0 sinon & \textbf{16.17} \\
\texttt{community\_size\_source} & Taille de la communauté du nœud source & 1.08 \\
\texttt{community\_size\_target} & Taille de la communauté du nœud cible & 1.20 \\
\bottomrule
\end{tabular}
\caption{Features de communauté (Louvain, 43 communautés détectées)}
\end{table}

\texttt{same\_community} est l'une des features les plus discriminantes (ratio de 16) : 79\% des paires positives sont dans la même communauté, contre seulement 5\% des négatives. C'est cohérent : les acteurs qui co-apparaissent sur Wikipedia forment naturellement des clusters thématiques.

\subsection{Famille 4 : Embeddings SVD (3 features)}

Pour capturer la structure \textbf{latente} du graphe, j'applique une décomposition SVD tronquée à la matrice d'adjacence. Cette technique projette les nœuds dans un espace de dimension réduite ($k = 64$) où la proximité reflète la similarité structurelle.

\begin{equation}
\mathbf{A} \approx \mathbf{U} \boldsymbol{\Sigma} \mathbf{V}^T, \quad \text{embedding de } u : \mathbf{e}_u = \mathbf{u}_i \sqrt{\sigma_i}
\end{equation}

\begin{table}[H]
\centering
\begin{tabular}{p{5cm}p{6cm}p{2cm}}
\toprule
\textbf{Feature} & \textbf{Formule} & \textbf{Ratio} \\
\midrule
\texttt{svd\_dot\_product} & $\mathbf{e}_u \cdot \mathbf{e}_v$ & \textbf{24.66} \\
\texttt{svd\_distance} & $\|\mathbf{e}_u - \mathbf{e}_v\|_2$ & 3.06 \\
\texttt{svd\_cosine\_similarity} & $\cos(\mathbf{e}_u, \mathbf{e}_v)$ & 3.84 \\
\bottomrule
\end{tabular}
\caption{Features d'embeddings SVD ($k=64$)}
\end{table}

Le \textbf{produit scalaire SVD} est la 2\textsuperscript{e} feature la plus discriminante (ratio 24.66). Les embeddings capturent des patterns de connectivité que les features locales (voisins communs) ne voient pas : deux acteurs sans voisin commun direct mais structurellement similaires auront des embeddings proches.

\subsection{Famille 5 : Centralité (9 features)}

Ces features caractérisent l'\textbf{importance} individuelle de chaque nœud dans le graphe.

\begin{table}[H]
\centering
\small
\begin{tabular}{p{5cm}p{6cm}p{2cm}}
\toprule
\textbf{Feature} & \textbf{Description} & \textbf{Ratio} \\
\midrule
\texttt{pagerank\_source/target} & Score PageRank de chaque nœud & 1.77 / 10.3 \\
\texttt{pagerank\_sum/diff/product} & Combinaisons & 6.3--13.6 \\
\texttt{clustering\_coeff\_*} & Coefficient de clustering local & $\sim$1.2 \\
\bottomrule
\end{tabular}
\caption{Features de centralité}
\end{table}

Le \textbf{PageRank} est particulièrement informatif : les paires positives impliquent des nœuds de plus haut PageRank. Le \textbf{coefficient de clustering} mesure à quel point les voisins d'un nœud sont interconnectés — un indicateur de densité locale du réseau.

% ============================================================
% 4. ANALYSE DES FEATURES
% ============================================================
\newpage
\section{Analyse des Features}

\subsection{Distributions par label}

La figure~\ref{fig:feat_dist} compare les distributions des features clés entre les paires avec lien (vert) et sans lien (rouge).

\begin{figure}[H]
\centering
\includegraphics[width=\textwidth]{figures/fig_04.png}
\caption{Distributions des 8 features les plus importantes. Les paires positives (vert) montrent des distributions nettement distinctes pour les features topologiques et de communauté, tandis que les features d'attributs (cosine similarity) sont moins discriminantes.}
\label{fig:feat_dist}
\end{figure}

\subsection{Matrice de corrélation}

\begin{figure}[H]
\centering
\includegraphics[width=0.85\textwidth]{figures/fig_05.png}
\caption{Matrice de corrélation entre les 36 features. On observe des blocs de corrélation au sein des familles (degrés, PageRank) mais une faible corrélation inter-familles, confirmant la complémentarité de l'approche.}
\label{fig:corr}
\end{figure}

\textbf{Observations clés} :
\begin{itemize}
    \item Les features de degré (\texttt{degree\_*}) sont fortement corrélées entre elles, ce qui est attendu.
    \item Les features de PageRank corrèlent avec les degrés — le PageRank est essentiellement une version normalisée du degré dans un graphe non dirigé.
    \item Les features d'attributs (cosine, Jaccard features) forment un bloc indépendant — elles apportent un signal orthogonal à la topologie.
    \item Les features SVD sont relativement indépendantes des autres, confirmant qu'elles capturent une information latente complémentaire.
\end{itemize}

\subsection{Contribution de chaque famille}

Pour quantifier l'apport de chaque famille, j'ai entraîné le modèle séparément sur chacune et mesuré l'AUC en validation croisée.

\begin{figure}[H]
\centering
\includegraphics[width=0.85\textwidth]{figures/fig_08.png}
\caption{AUC de chaque famille de features utilisée seule. La topologie est la plus performante, complétée par la communauté et les embeddings SVD.}
\label{fig:ablation}
\end{figure}

La \textbf{topologie du graphe} est de loin la famille la plus puissante, mais les autres familles apportent une amélioration marginale qui peut faire la différence en compétition. Cette étude d'ablation confirme que toutes les familles contribuent et méritent d'être conservées.

% ============================================================
% 5. MODÉLISATION
% ============================================================
\newpage
\section{Modélisation}

\subsection{Modèles comparés}

J'ai évalué 4 algorithmes de machine learning classique :

\begin{enumerate}
    \item \textbf{Régression Logistique} : modèle linéaire, baseline rapide. Paramètres : $C=1.0$, solver LBFGS.
    \item \textbf{Random Forest} : ensemble de 500 arbres de décision. Paramètres : profondeur max 12, min 5 échantillons par feuille.
    \item \textbf{XGBoost} : gradient boosting optimisé. Paramètres : 500 arbres, profondeur 6, learning rate 0.1, subsample 80\%.
    \item \textbf{LightGBM} : gradient boosting rapide de Microsoft. Paramètres : 500 arbres, profondeur 8, learning rate 0.1, subsample 80\%.
\end{enumerate}

\subsection{Validation croisée}

La validation est réalisée avec un \textbf{StratifiedKFold à 5 folds} (seed=42) pour maintenir l'équilibre des classes dans chaque fold. La métrique reportée est l'AUC-ROC.

\begin{table}[H]
\centering
\begin{tabular}{lccc}
\toprule
\textbf{Modèle} & \textbf{AUC moyenne} & \textbf{Écart-type} & \textbf{Temps} \\
\midrule
\textbf{LightGBM} & \textbf{0.9982} & 0.0006 & 6.6s \\
XGBoost & 0.9980 & 0.0005 & 7.9s \\
Random Forest & 0.9935 & 0.0015 & 12.8s \\
Régression Logistique & 0.9852 & 0.0027 & 0.6s \\
\bottomrule
\end{tabular}
\caption{Résultats de la validation croisée 5-fold}
\end{table}

Tous les modèles obtiennent une excellente AUC, mais \textbf{LightGBM} et \textbf{XGBoost} sont en tête. La différence entre la CV (0.998) et le score Kaggle (0.82) s'explique par le fait que certaines features globales (SVD, communauté) sont calculées sur le graphe complet, ce qui donne un léger avantage en CV qui ne se retrouve pas intégralement sur les données de test Kaggle, où les arêtes manquantes sont différentes.

\subsection{Courbes ROC}

\begin{figure}[H]
\centering
\includegraphics[width=0.8\textwidth]{figures/fig_06.png}
\caption{Courbes ROC des 4 modèles. Tous sont très proches du coin supérieur gauche (classification parfaite), avec les modèles de boosting légèrement en avance.}
\label{fig:roc}
\end{figure}

\subsection{Importance des features}

Le modèle retenu (\textbf{LightGBM}) fournit une mesure d'importance pour chaque feature, basée sur le nombre de fois qu'elle est utilisée comme critère de split dans les arbres.

\begin{figure}[H]
\centering
\includegraphics[width=0.75\textwidth]{figures/fig_07.png}
\caption{Importance des 36 features dans le modèle LightGBM final. Les features topologiques (preferential attachment, degrés) et de communauté (same\_community) dominent.}
\label{fig:importance}
\end{figure}

\textbf{Top 5 features} : \texttt{preferential\_attachment}, \texttt{same\_community}, \texttt{degree\_max}, \texttt{degree\_target}, \texttt{svd\_dot\_product}. Ce classement confirme que la structure du graphe est le signal dominant, avec un apport important de la détection de communautés et des embeddings SVD.

% ============================================================
% 6. GÉNÉRATION DES PRÉDICTIONS
% ============================================================
\newpage
\section{Génération des Prédictions}

\subsection{Processus}

\begin{enumerate}
    \item Le modèle LightGBM est entraîné sur \textbf{l'intégralité} des données d'entraînement (10\,496 paires).
    \item Les features sont extraites pour les 3\,498 paires de test en mode \textbf{inférence} (sans retrait d'arêtes, puisque ces paires ne sont pas dans le graphe).
    \item Le modèle prédit des \textbf{probabilités} $P(\text{lien} | \text{features})$ pour chaque paire.
    \item Le fichier de soumission est généré au format : \texttt{ID, Predicted}.
\end{enumerate}

\begin{figure}[H]
\centering
\includegraphics[width=0.8\textwidth]{figures/fig_09.png}
\caption{Distribution des probabilités prédites sur les données de test. On observe une séparation bimodale avec un pic autour de 0 (non-liens) et un pic autour de 1 (liens).}
\label{fig:predictions}
\end{figure}

La distribution bimodale confirme que le modèle est \textbf{confiant} dans ses prédictions : la plupart des paires sont prédites avec une probabilité très proche de 0 ou de 1, avec peu de cas ambigus.

\subsection{Résultat Kaggle}

Le score obtenu sur le leaderboard Kaggle est : $\boxed{\text{AUC} = 0.82}$

% ============================================================
% 7. ARCHITECTURE DU CODE
% ============================================================
\section{Architecture du Code}

Le projet est organisé de manière modulaire :

\begin{table}[H]
\centering
\small
\begin{tabular}{p{5cm}p{8cm}}
\toprule
\textbf{Fichier} & \textbf{Rôle} \\
\midrule
\texttt{src/data\_loader.py} & Chargement des données, construction du graphe NetworkX \\
\texttt{src/feature\_extractor.py} & Extraction des 36 features (5 familles), gestion du leakage \\
\texttt{src/train.py} & Entraînement, CV 5-fold, comparaison de 4 modèles \\
\texttt{src/generate\_submission.py} & Génération du CSV de soumission Kaggle \\
\texttt{src/create\_notebook.py} & Création programmatique du notebook \\
\texttt{notebook.ipynb} & Notebook EDA + analyse + résultats \\
\bottomrule
\end{tabular}
\caption{Structure du code source}
\end{table}

Exécution :
\begin{lstlisting}[language=bash]
venv\Scripts\python.exe src\train.py              # Entrainer
venv\Scripts\python.exe src\generate_submission.py # Soumettre
\end{lstlisting}

% ============================================================
% 8. CONCLUSION
% ============================================================
\section{Conclusion et Pistes d'Amélioration}

\subsection{Bilan}

Cette approche basée sur un \textbf{feature engineering extensif} et du \textbf{gradient boosting} (LightGBM) obtient un score Kaggle de \textbf{0.82 AUC}. Les principaux facteurs de succès sont :

\begin{enumerate}
    \item La \textbf{combinaison de 5 familles de features} complémentaires, allant de la topologie locale aux embeddings globaux.
    \item L'utilisation d'\textbf{heuristiques classiques} de link prediction (Adamic-Adar, Resource Allocation) qui exploitent efficacement la structure du graphe sparse.
    \item La \textbf{détection de communautés} via Louvain, qui capture la structure mésoscopique du réseau.
    \item Les \textbf{embeddings SVD}, qui projettent la structure globale du graphe dans un espace latent informatif.
\end{enumerate}

\subsection{Pistes d'amélioration}

\begin{itemize}
    \item \textbf{Ensemble de modèles} : combiner LightGBM, XGBoost et Random Forest par stacking ou voting pour réduire la variance.
    \item \textbf{Hyperparameter tuning} : utiliser Optuna ou une recherche bayésienne pour optimiser les paramètres du modèle.
    \item \textbf{Features supplémentaires} : voisins communs à 2 sauts, centralité de betweenness, spectral features, embeddings Node2Vec.
    \item \textbf{Validation croisée améliorée} : recalculer les features SVD et communauté à chaque fold en excluant les arêtes de validation, pour un score CV plus réaliste.
    \item \textbf{Sélection de features} : supprimer les features redondantes (degrés corrélés, PageRank $\approx$ degré) pour réduire le bruit.
\end{itemize}

\end{document}
